\lecture{22}{Wed 05 Nov 2025 12:20}{Covector fields}

Recall: given a path $\gamma : [a,b] \to M$ in a manifolds and a covector field $\omega \in \mathfrak{X} ^*(M)$, we can define the \emph{line integral} of $\omega $ over $\gamma $ by
\[
	\int_\gamma \omega =\int_{a} ^b \gamma ^*\omega 
.\]
In vector calculus, the line integral of a vector field $V : \mathbb{R} ^m \to \mathbb{R} ^n$ over a path $\gamma : [a,b] \to \mathbb{R} ^n$ was given by
\[
	\int_\gamma  V\cdot  \,\mathrm{d} \mathbf{s} =\int_a^b V(\gamma (t)) \frac{\partial \gamma }{\partial t}  \,\mathrm{d} t
.\]
It would be interesting if our formula basically became the same. The reason vector calc used vectors instead of covectors was because we had a nice inner product. On manifolds, we don't, so we have to use covectors and do more work.

Given a vector space $V$, we can define a bilinear map $g : V\times V \to \mathbb{R} $ by $g(v_1,v_2)=v_1\cdot v_2$ the dot product. We can then define a map $g^\flat : V \to V^*$ by $v\mapsto g(v,-)$. This is linear and has trivial kernel, so its an isomorphism of vector spaces.

Let's move on to considering conservative fields. In calculus 3, a vector field $V$ was called conservative if it had a potential function $V=\nabla f$ for some $f$, and it was path independent if and only if it was conservative.

\begin{proposition}\label{3.1.5}
	If $\gamma : [a,b] \to M$ is piecewise smooth and $\omega $ is a covector field on $M$, then
	\[
		\int_\gamma \omega =\int_a^b \omega_{\gamma (t)} (\gamma '(t)) \,\mathrm{d} t
	.\]
	This reduces this to just a calculus 1 integral.
\end{proposition}

\begin{proposition}\label{3.1.6}
	Let $M$ be a smooth manifold, $\omega \in \mathfrak{X} ^*(M)$, and $\gamma : [a,b] \to M$ piecewise smooth. Let $\widetilde{\gamma } $ be a \emph{forwards representation} of $\gamma $. Then
	\[
		\int_{\widetilde{\gamma } } \omega =\int_\gamma \omega 
	.\]
	If $\widetilde{\gamma } $ is a \emph{backwardsd reparameterization}, then
	\[
		\int_{\widetilde{\gamma } } \omega =-\int_\gamma \omega
	.\]
\end{proposition}
\begin{definition}\label{3.1.7}
	Given $\gamma : [a,b] \to M$ a piecewise smooth curve, a forwards reparameterization is a reparameterization which is monotonic, so if $\gamma$ is in terms of $t$ and we change $t\to s$, then as $t$ increases, $s$ increases, and vice versa. A backwards reparameterization is the opposite.
\end{definition}

\begin{theorem}\label{3.1.8}
	Let $M$ be a smooth manifold, $\gamma : [a,b] \to M$ piecewise smooth, and $f : M \to \mathbb{R} $ a smooth function. Then recalling that $\mathrm{d} f\in \mathfrak{X} ^*(M)$,
	\[
		\int_\gamma \mathrm{d} f = (f\circ \gamma )(b)-(f\circ \gamma )(a)
	.\]
\end{theorem}
\begin{proof}
	By our definition of d and \cref{3.1.5},
	\[
		\int_\gamma  \mathrm{d} f = \int_a^b \mathrm{d} f_{\gamma (t)} (\gamma '(t)) \,\mathrm{d} t = \int_a^b (f\circ \gamma )'(t) \,\mathrm{d} t
	.\]
	Applying the fundamentl theorem of calculus recovers the statement.
\end{proof}
\begin{corollary}\label{3.1.9}
	If $\gamma (a)=\gamma (b)$, then
	\[
		\int_\gamma \mathrm{d} f =0
	.\]
\end{corollary}

\begin{definition}\label{3.1.10}
	Let $\omega \in \mathfrak{X} ^*(M)$. Then $\omega $ is \emph{exact} if and only if $\omega =\mathrm{d} f$ for some $f\in C^\infty (M)$. We say $f$ is the \emph{potential function} of $\omega $. We say $\omega $ is \emph{conservative} if and only if
	\[
		\int_\gamma  \omega =0
	\]
	for all piecewise smooth \emph{closed} curves (loops) $\gamma $.
\end{definition}

In calculus, exactness and conservativeness were identical. On manifolds, we have aleady shown that exact $\implies $ conservative. Is the converse true?

\begin{proposition}\label{3.1.11}
	A covector field $\omega \in \mathfrak{X} ^*(M)$ is conservative if and only if
	\[
		\int_\gamma \omega =\int_{\widetilde{\gamma } } \omega 
	\]
	for all $\gamma ,\widetilde{\gamma } $ starting and ending at the same points.
\end{proposition}
\begin{proof}
	Using \cref{3.1.3} and \cref{3.1.6}, form a closed loop be reparameterizing one of them backwards and adding them together.
\end{proof}

\begin{theorem}\label{3.1.12}
	A smooth covector field $\omega $ on $M$ is conservative if and only if it is exact.
\end{theorem}
\begin{proof}
	We already have one direction, so let $\omega $ be conservative. We can assume $M$ is connected by taking a connected component. Let $\gamma : [a,b] \to M$ be piecewise smooth, and write
	\[
		\int_{\gamma (a)}^{\gamma (b)} \omega =\int_\gamma \omega
	.\]
	Since $\omega $ is conservative, this notation is well-defined. Pick a basepoint $p_0\in M$, and define
	\[
		f(q)\coloneqq \int_{p_0} ^q\omega
	.\]
	We claim that $\mathrm{d} f=\omega $. Choose some $q_0\in M$, and choose a chart $(U,\varphi )$ of $q_0$. Write $x^i$ for the coordinates. The blackboard got messy and I got lost but it follows. Basically you write it in coordinates then use the same argument we did in calculus 3 to show the proof. I don't remember the calculus three argument...
\end{proof}

